\documentclass[11pt,a4paper]{article}
\usepackage[utf8]{inputenc}
\usepackage[T1]{fontenc}
\usepackage{fouriernc}
\usepackage{amsfonts}
\usepackage[french]{babel}
\frenchbsetup{StandardLists=true}
\usepackage{enumitem}
\usepackage{amssymb}
\usepackage{amsmath}
\usepackage[left=2cm,right=2cm,top=2cm,bottom=2cm]{geometry}
\usepackage{multirow}
\usepackage[dvipsnames]{xcolor}
\usepackage{graphicx}
\usepackage{setspace}
\singlespacing
\usepackage{fancyhdr}
\setlength{\headheight}{14pt}
\usepackage{tcolorbox}
\usepackage{cancel}

\pagestyle{fancy}
\renewcommand\headrulewidth{1pt}
\fancyhead[L]{Lycée Denis Diderot}
\fancyhead[R]{Classe 2BTS}
\fancyfoot[C]{Corrigé - Page \thepage}
\fancyfoot[R]{Année 2025-2026}
\fancyfoot[L]{Alexis RUIZ}

\setlength{\parindent}{0cm}
\renewcommand{\arraystretch}{1.8}
\DeclareMathOperator{\e}{e}

% Boîte pour encadrer les résultats
\newcommand{\reponse}[1]{%
  \begin{tcolorbox}[colback=blue!5!white, colframe=blue!50!black,
    boxrule=0.5pt, left=4pt, right=4pt, top=2pt, bottom=2pt]
  #1
  \end{tcolorbox}%
}

\begin{document}

\bigskip
\begin{center}
  \LARGE\textbf{Corrigé – Équations différentielles (2 h)}\\[0.3cm]
\end{center}
\normalsize
\hrule
\vspace{0.3cm}

%%%%%%%%%%%%%%%%%%%%%%%%%%%%%%%%%%%%%%%%%%%%%%%%%%
\fcolorbox[gray]{0}{0.9}{\textbf{Exercice 1}} \quad
$y' + 2y = -5\e^{-2x}$

\medskip

\textbf{1. Résolution de (E$_0$) : $y' + 2y = 0$}

L'équation caractéristique est $r + 2 = 0$, soit $r = -2$.

\reponse{Les solutions de (E$_0$) sont : $y_0 = C\e^{-2x}$, \quad $C \in \mathbb{R}$}

\bigskip

\textbf{2. Vérification que $g(x) = -5x\e^{-2x}$ est solution particulière de (E)}

On calcule $g'(x)$ par la formule de dérivation d'un produit :
\[
  g'(x) = -5\e^{-2x} + (-5x)\times(-2)\e^{-2x} = -5\e^{-2x} + 10x\e^{-2x}
\]

On substitue dans (E) :
\[
  g'(x) + 2g(x) = \bigl(-5\e^{-2x} + 10x\e^{-2x}\bigr) + 2\bigl(-5x\e^{-2x}\bigr)
  = -5\e^{-2x} + \cancel{10x\e^{-2x}} - \cancel{10x\e^{-2x}}
  = -5\e^{-2x}
\]

On obtient bien $-5\e^{-2x}$, donc $g$ est une solution particulière de (E). \hfill $\square$

\bigskip

\textbf{3. Ensemble des solutions de (E)}

La solution générale de (E) est la somme de la solution générale de (E$_0$) et d'une solution particulière :

\reponse{$y = y_0 + g = C\e^{-2x} - 5x\e^{-2x} = (C - 5x)\e^{-2x}$, \quad $C \in \mathbb{R}$}

\bigskip

\textbf{4. Solution particulière vérifiant $f(0) = 1$}

On pose $f(x) = (C - 5x)\e^{-2x}$ et on applique la condition initiale :
\[
  f(0) = (C - 0)\cdot\e^{0} = C = 1
\]

\reponse{$f(x) = (1 - 5x)\e^{-2x}$}

\newpage

%%%%%%%%%%%%%%%%%%%%%%%%%%%%%%%%%%%%%%%%%%%%%%%%%%
\fcolorbox[gray]{0}{0.9}{\textbf{Exercice 2}} \quad
$y' - 2y = -x^{2}$

\medskip

\textbf{1. Résolution de (E$_0$) : $y' - 2y = 0$}

L'équation caractéristique est $r - 2 = 0$, soit $r = 2$.

\reponse{Les solutions de (E$_0$) sont : $y_0 = C\e^{2x}$, \quad $C \in \mathbb{R}$}

\bigskip

\textbf{2. Recherche d'une solution particulière $y_p(x) = Ax^2 + Bx + C$}

On calcule $y_p'(x) = 2Ax + B$. On substitue dans (E) :
\[
  y_p' - 2y_p = (2Ax + B) - 2(Ax^2 + Bx + C)
  = -2Ax^2 + (2A - 2B)x + (B - 2C)
\]

On identifie terme à terme avec $-x^2 + 0 \cdot x + 0$ :
\[
  \begin{cases}
    -2A = -1 \\
    2A - 2B = 0 \\
    B - 2C = 0
  \end{cases}
  \implies
  \begin{cases}
    A = \dfrac{1}{2} \\[4pt]
    B = A = \dfrac{1}{2} \\[4pt]
    C = \dfrac{B}{2} = \dfrac{1}{4}
  \end{cases}
\]

\reponse{$y_p(x) = \dfrac{1}{2}x^2 + \dfrac{1}{2}x + \dfrac{1}{4}$}

\bigskip

\textbf{3. Ensemble des solutions de (E)}

\reponse{$y = C\e^{2x} + \dfrac{1}{2}x^2 + \dfrac{1}{2}x + \dfrac{1}{4}$, \quad $C \in \mathbb{R}$}

\bigskip

\textbf{4. Solution particulière vérifiant $f(0) = 1$}

On applique la condition initiale :
\[
  f(0) = C\cdot\e^{0} + 0 + 0 + \frac{1}{4} = C + \frac{1}{4} = 1 \implies C = \frac{3}{4}
\]

\reponse{$f(x) = \dfrac{3}{4}\e^{2x} + \dfrac{1}{2}x^2 + \dfrac{1}{2}x + \dfrac{1}{4}$}

\newpage

%%%%%%%%%%%%%%%%%%%%%%%%%%%%%%%%%%%%%%%%%%%%%%%%%%
\fcolorbox[gray]{0}{0.9}{\textbf{Exercice 3}} \quad
$y' - y = (3x-1)\e^{x}$

\medskip

\textbf{1. Résolution de (E$_0$) : $y' - y = 0$}

L'équation caractéristique est $r - 1 = 0$, soit $r = 1$.

\reponse{Les solutions de (E$_0$) sont : $y_0 = C\e^{x}$, \quad $C \in \mathbb{R}$}

\bigskip

\textbf{2. Recherche d'une solution particulière $y_p(x) = (Ax^2 + Bx + C)\e^{x}$}

On calcule $y_p'(x)$ par la formule de dérivation d'un produit :
\begin{align*}
  y_p'(x) &= (2Ax + B)\e^x + (Ax^2 + Bx + C)\e^x \\
           &= \bigl(Ax^2 + (2A+B)x + (B+C)\bigr)\e^x
\end{align*}

On substitue dans (E) :
\begin{align*}
  y_p' - y_p &= \bigl(Ax^2 + (2A+B)x + (B+C)\bigr)\e^x - (Ax^2 + Bx + C)\e^x \\
              &= \bigl(\cancel{Ax^2} + (2A+\cancel{B})x + (\cancel{B}+\cancel{C})\bigr)\e^x
               - (\cancel{Ax^2} + \cancel{Bx} + \cancel{C})\e^x \\
              &= (2Ax + B)\e^x
\end{align*}

On identifie terme à terme avec $(3x - 1)\e^x$ :
\[
  \begin{cases}
    2A = 3 \\
    B = -1
  \end{cases}
  \implies
  \begin{cases}
    A = \dfrac{3}{2} \\[4pt]
    B = -1
  \end{cases}
  \quad (C \text{ est libre, on choisit } C = 0)
\]

\reponse{$y_p(x) = \left(\dfrac{3}{2}x^2 - x\right)\e^{x}$}

\bigskip

\textbf{3. Ensemble des solutions de (E)}

\reponse{$y = \left(C + \dfrac{3}{2}x^2 - x\right)\e^{x}$, \quad $C \in \mathbb{R}$}

\bigskip

\textbf{4. Solution particulière vérifiant $f(0) = -1$}

On applique la condition initiale :
\[
  f(0) = \left(C + 0 - 0\right)\e^{0} = C = -1
\]

\reponse{$f(x) = \left(\dfrac{3}{2}x^2 - x - 1\right)\e^{x}$}

\end{document}
